\section{Theoretical Analysis}
\label{sec:theory}
\subsection{Wasserstein bound}
Because SGLD targets a continuous distribution, we need a metric that quantifies how close the distribution $\nu_k$ of the $k$-th iterate is to the true posterior $\pi$. The 2-Wasserstein distance $W_2(\nu_k,\pi)$ is a natural choice: it metrizes weak convergence while accounting for the geometry of the parameter space, and its squared form admits a clean coupling argument that translates the oracle-noise variance of our gradient estimator directly into a recursive bound on sampling error. Controlling $W_2$ therefore lets us compare SGLD-CV2 against Baker's SGLD-CV on equal footing and makes the benefit of the quartic variance bound (\Cref{prop:quartic}) precise in terms of the distribution's actual distance to stationarity.
From lectures, our sampling target is:
$$
    \pi(\theta) \propto e^{-U(\theta)}, \quad \theta \in \mathbb{R}^d,
$$
where $U : \mathbb{R}^d \to \mathbb{R}$ is differentiable.
Recall the overdamped Langevin diffusion:
$$
    d\theta_t = -\nabla U(\theta_t)\,dt + \sqrt{2}\,dB_t.
$$
In the integral form:
$$
    \theta_t = \theta_0 - \int_0^t \nabla U(\theta_s)\,ds + \sqrt{2}\,B_t.
$$
From \citet{welling2011}:
$$U(\theta) = U_0(\theta) + \sum_{i=1}^N U_i(\theta)$$
where $U_0(\theta) = -\log p(\theta)$ is the prior term and $U_i(\theta) = -\log p(y_i \mid \theta)$ are the likelihood terms. And so:
\begin{align}
    \nabla U(\theta) &= \nabla U_0(\theta) + \sum_{i=1}^N \nabla U_i(\theta), \\
    \nabla^2 U(\theta) &= \nabla^2 U_0(\theta) + \sum_{i=1}^N \nabla^2 U_i(\theta).
\end{align}
At each iteration, SGLD uses the mini-batch gradient
\begin{equation}
    \nabla \tilde{U}(\theta_k) = \nabla U_0(\theta_k) 
    + \frac{N}{n}\sum_{i \in S_k} \nabla U_i(\theta_k)
\end{equation}
where $S_k \subset \{1,\ldots,N\}$ is a uniformly sampled mini-batch of size 
$|S_k| = n$. This is an unbiased estimator of the full-batch gradient, i.e. $\mathbb{E}_{S_k}\!\left[\nabla \hat{U}(\theta_k)\right] = \nabla U(\theta_k)$.

Our proposed \textbf{second-order control variate} is:
$$C_k = \frac{N}{n}\sum_{i \in S_k} [\nabla U_i(\hat\theta) + \nabla^2 U_i(\hat\theta)(\theta_k - \hat\theta)]$$
with deterministic expectation 
$$\mathbb{E}[C_k] = \left[\nabla U(\hat\theta) - \nabla U_0(\hat\theta)\right] + \left[\nabla^2 U(\hat\theta) - \nabla^2 U_0(\hat\theta)\right](\theta_k - \hat\theta)$$
Applying $\nabla\hat{U}_{\text{CV2}}(\theta_k) = \nabla\tilde{U}(\theta_k) - C_k + \mathbb{E}[C_k]$ as a control-variated unbiased estimator of the full-batch gradient:
\begin{align}
\nabla\hat{U}_{\text{CV2}}(\theta_k) 
    &= \nabla U_0(\theta_k) + \bigl[\nabla U(\hat\theta) 
        + \nabla^2 U(\hat\theta)(\theta_k - \hat\theta)\bigr] - \bigl[\nabla U_0(\hat\theta) 
        + \nabla^2 U_0(\hat\theta)(\theta_k - \hat\theta)\bigr] \notag\\
    &\quad + \frac{N}{n}\sum_{i \in S_k}\bigl(\nabla U_i(\theta_k) 
        - \nabla U_i(\hat\theta) 
        - \nabla^2 U_i(\hat\theta)(\theta_k - \hat\theta)\bigr).
    \label{eq:cv2-theory}
\end{align}
% \jay{I think we have duplicate expositions above on SGLD-CV2 with section 3. Michelle since you put this in first, feel free to trim mine.}
\begin{assumption}\label{ass:convex}[Strongly log-concave target, i.e.\ strongly convex potential]
There exist positive constants $m$ and $M$ such that the following holds for the potential $U$:
\begin{align}
    U(\theta) - U(\theta') - \nabla U(\theta')^\top (\theta - \theta') &\geq \frac{m}{2} \|\theta - \theta'\|^2 \\
    \|\nabla U(\theta) - \nabla U(\theta')\| &\leq M \|\theta - \theta'\|.
\end{align}
for all $\theta, \theta' \in \mathbb{R}^d$.
\end{assumption}
We further need a Lipschitz condition for each of the likelihood terms in order to bound the variance of our control-variate estimator of the gradient.
\begin{assumption}[Lipschitz Gradients]
\label{ass:grad}
    There exist constants $L_0, \ldots, L_N$ such that
\begin{equation}
    \|\nabla U_i(\theta) - \nabla U_i(\theta')\| \leq L_i \|\theta - \theta'\|, 
    \quad \text{for } i = 0, \ldots, N.
\end{equation}
\end{assumption}
\begin{assumption}[Lipschitz Hessians]
\label{ass:hess}
    There exist constants $M_0, \ldots, M_N$ such that
\begin{equation}
|\nabla^2 U_i(\theta) - \nabla^2 U_i(\theta')| \leq M_i |\theta - \theta'|, \quad \text{for } i = 0, \ldots, N.
\end{equation}
\end{assumption}

Using \Cref{ass:hess} we are able to derive a bound on the variance of the gradient estimate of SGLD-CV. This bound is formally stated in Proposition 1.

\begin{proposition}\label{prop:quartic}
\textit{Under Assumption 3. Let $\theta_k$ be the state of SGLD-CV2 at the $k^{th}$ 
iteration, with stepsize $h$ and centering value $\hat\theta$. Assume we estimate 
the gradient using the second-order control variate estimator. Define $\xi_k := \nabla \hat{U}(\theta_k) - \nabla U(\theta_k)$, 
so that $\xi_k$ measures the noise in the gradient estimate $\nabla \hat{U}$ 
and has mean $0$. Then for all $\theta_k, \hat{\theta} \in \mathbb{R}^d$, for all 
$k = 1, \ldots, K$ we have}
\begin{equation}\label{eq:lemma1}
\mathbb{E}\|\xi_k\|^2 \leq \frac{N}{4n} \left( \sum_{i=1}^N M_i^2 \right) \mathbb{E}\|\theta_k - \hat\theta\|^4.
\end{equation}
\end{proposition}

\begin{proof}
Our proof follows similarly to \citet{baker2017}.
Note that:
$$\nabla^2 U(\hat\theta)(\theta_k - \hat\theta) = \sum_{i=1}^N \nabla^2 U_i(\hat\theta)(\theta_k - \hat\theta)\bigr)$$
\begin{align*}
\mathbb{E}\|\xi_k\|^2 
&= \mathbb{E}\left\|\nabla\hat{U}_{\text{CV2}}(\theta_k) - \nabla U(\theta_k)\right\|^2 \\
&= \mathbb{E}\Bigg\|
    \nabla U_0(\theta_k) + \bigl[\nabla U(\hat\theta) 
        + \nabla^2 U(\hat\theta)(\theta_k - \hat\theta)\bigr] - \bigl[\nabla U_0(\hat\theta) 
        + \nabla^2 U_0(\hat\theta)(\theta_k - \hat\theta)\bigr] \notag\\
    &\quad + \frac{N}{n}\sum_{i \in S_k}\bigl(\nabla U_i(\theta_k) 
        - \nabla U_i(\hat\theta) 
        - \nabla^2 U_i(\hat\theta)(\theta_k - \hat\theta)\bigr) - \nabla U(\theta_k)
    \Bigg\|^2 \\
&= \mathbb{E}\Bigg\|
    \bigl[\nabla U_0(\theta_k) - \nabla U_0(\hat\theta) - \nabla^2 U(\hat\theta)(\theta_k - \hat\theta)\bigr] \\
&\quad + \frac{1}{n}\sum_{i \in S_k} N\bigl[\nabla U_i(\theta_k) 
        - \nabla U_i(\hat\theta) 
        - \nabla^2 U_i(\hat\theta)(\theta_k - \hat\theta)\bigr] \notag\\
    &\quad - \bigl[\nabla U(\theta_k) - \nabla U(\hat\theta) - \nabla^2 U_0(\hat\theta)(\theta_k - \hat\theta)\bigr]
    \Bigg\|^2 \\
\end{align*}
Applying inequality $\mathbb{E}\|\frac{1}{n}\sum Z_i - \mathbb{E}[Z]\|^2 \leq \frac{1}{n^2}\mathbb{E}\sum\|Z_i - \mathbb{E}[Z]\|^2$:

\nopagebreak[4]
\begin{align*}
\begin{aligned}
\mathbb{E}\|\xi_k\|^2 &\leq \frac{1}{n^2} \mathbb{E} \sum_{i \in S_k} \Bigg\| \Bigl[\nabla U(\theta_k) - \nabla U(\hat\theta) - \nabla^2 U(\hat\theta)(\theta_k - \hat\theta)\Bigr] \\
&\quad - \biggl( \Bigl[\nabla U_0(\theta_k) - \nabla U_0(\hat\theta) - \nabla^2 U_0(\hat\theta)(\theta_k - \hat\theta)\Bigr] \\ &\quad \quad + N\Bigl[\nabla U_i(\theta_k) - \nabla U_i(\hat\theta) - \nabla^2 U_i(\hat\theta)(\theta_k - \hat\theta)\Bigr] \biggr) \Bigg\|^2.
\end{aligned}
\end{align*}
For a uniformly sampled data point $I$:
$$
    R = N \Bigl[ \nabla U_I(\theta_k) - \nabla U_I(\hat\theta) 
    - \nabla^2 U_I(\hat\theta)(\theta_k - \hat\theta) \Bigr].
$$
Taking the expectation over the uniform distribution of $I$:
\begin{align*}
    \mathbb{E}_I[R] 
    &= \sum_{i=1}^N \frac{1}{N} \Bigg( N \Bigl[ \nabla U_i(\theta_k) 
        - \nabla U_i(\hat\theta) 
        - \nabla^2 U_i(\hat\theta)(\theta_k - \hat\theta) \Bigr] \Bigg) \\
    &= \sum_{i=1}^N \Bigl[ \nabla U_i(\theta_k) 
        - \nabla U_i(\hat\theta) 
        - \nabla^2 U_i(\hat\theta)(\theta_k - \hat\theta) \Bigr].
\end{align*}
Since $U = U_0 + \sum_{i=1}^N U_i$:
\begin{align*}
    \mathbb{E}_I[R] 
    &= \Bigl[ \nabla U(\theta_k) - \nabla U(\hat\theta) 
        - \nabla^2 U(\hat\theta)(\theta_k - \hat\theta) 
        \Bigr] \\
&\quad - \Bigl[ \nabla U_0(\theta_k) - \nabla U_0(\hat\theta) 
        - \nabla^2 U_0(\hat\theta)(\theta_k - \hat\theta) 
        \Bigr].
\end{align*}
Substituting into the inequality, and using the fact that for any random variable $R$, 
we have that $\mathbb{E}\|R - \mathbb{E}R\|^2 \leq \mathbb{E}\|R\|^2$:
\begin{align*}
\mathbb{E}\|\xi_k\|^2 
&\leq \frac{1}{n^2}\,\mathbb{E}\sum_{i \in S_k}
    \left\|\mathbb{E}_I[R] - R_i\right\|^2 = \frac{1}{n^2}\,\mathbb{E}\sum_{i \in S_k}
    \left\|R_i - \mathbb{E}_I[R]\right\|^2 \\
&\leq \frac{1}{n^2}\,\mathbb{E}\sum_{i \in S_k}
    \left\|R_i\right\|^2 = \frac{1}{n}\,\mathbb{E}_I\left\|R_I\right\|^2
\end{align*}
 

Using this where $\mathbb{E}_I$ refers to 
expectation with respect to the sampled datum index, $I$ and and the Lipschitz results of Assumption 3, i.e. $$\left\| \nabla U_i(x) - \nabla U_i(y) - \nabla^2 U_i(y)(x - y) \right\| \leq \frac{M_i}{2} \|x - y\|^2$$
We have:
\begin{align}
\mathbb{E}\|\xi_k\|^2 
&\leq \frac{N}{n} \sum_{i=1}^N \left| \nabla U_i(\theta_k) - \nabla U_i(\hat\theta) - \nabla^2 U_i(\hat\theta)(\theta_k - \hat\theta) \right|^2 \notag\\
&\leq \frac{N}{n} \sum_{i=1}^N \left( \frac{M_i}{2} |\theta_k - \hat\theta|^2 \right)^2 \notag\\
&= \frac{N}{4n} \left( \sum_{i=1}^N M_i^2 \right) \|\theta_k - \hat\theta\|^4
\end{align}
from which the claimed bound follows from taking expectation of both sides.
\end{proof}

If Assumption 1 also holds, then we 
can choose $M = \sum_{i=0}^N L_i$, and our bound \eqref{eq:lemma1} implies
$$
    \mathbb{E}\|\xi_k\|^2 \leq \frac{M^2}{n}\,\mathbb{E}\left\|\theta_k - \hat{\theta}\right\|^2.
$$
We will use this form of the bound for the rest of the analysis. In order to consider how SGLD-CV2 scales with $N$ we need to make assumptions 
on the properties of the target and how these change with $N$. As in \citet{baker2017} we will focus on the following, strong, assumption that 
each likelihood term in the target is strongly log-concave. As we discuss 
later, our results apply under weaker conditions.

\begin{assumption}[Uniform Lipschitz Hessians]
\label{ass:hess-unif} \textit{Assume there exists a positive constant $H$ such that the Hessian Lipschitz constants for all likelihood terms are uniformly bounded:}
\begin{equation}
M_i \leq H, \quad \text{for all } i = 1, \ldots, N.
\end{equation}
\end{assumption}
Under this assumption, we can simplify the variance bound from Proposition 1. Since $\sum_{i=1}^N M_i^2 \leq N H^2$, our bound implies:
\begin{equation}
\mathbb{E}\|\xi_k\|^2 \leq \frac{N^2 H^2}{4n} \mathbb{E}\|\theta_k - \hat\theta\|^4.
\end{equation}
To see the theoretical benefit of using second-order control variates, we can compare this variance bound with the variance of the simple SGLD estimator and the first-order SGLD-CV estimator. If we assume that $\|\nabla U_i(\theta)\|$is bounded by some constant $\sigma$, Dubey et al. (2016) show that for standard SGLD:
\begin{equation}
\mathbb{E}\left\|\nabla\hat{U}(\theta) - \nabla U(\theta)\right\|^2 \leq \frac{2N^2\sigma^2}{n} \in O(N^2/n).
\end{equation}
We can see that the bound on the gradient estimate variance for SGLD-CV2 depends on the fourth power of the distance between $\theta_k$ and $\hat\theta$. Appealing to the Bernstein-von Mises theorem (Le Cam, 2012), under standard asymptotics we would expect the squared distance $\mathbb{E}\|\theta_k - \hat\theta\|^2$ to be $O(1/N)$, if $\hat\theta$ is within $O(N^{-1/2})$ of the posterior mean, once the MCMC algorithm has burnt in. Consequently, the fourth moment $\mathbb{E}\|\theta_k - \hat\theta\|^4$ becomes $O(N^{-2})$. Substituting this into our bound, we find:
\begin{equation}
\mathbb{E}\|\xi_k\|^2 \leq \frac{N^2 H^2}{4n} O(N^{-2}) = O(1/n).
\end{equation}
This suggests that using a second-order control variate gives an $O(N^2)$ reduction in variance compared to standard SGLD, and an $O(N)$ reduction compared to the first-order SGLD-CV, and this is crucial in the computational cost improvements we show in
the next section.
\subsection{Computational cost of SGLD-CV2}
We investigate the Wasserstein given by the SGLD-CV2 algorithm at each iteration and the true posterior as $N$ changes.

\textbf{Theorem 1} \textit{Under Assumption 1, suppose $h < 2m/(2M^2 + m^2)$. Let $\nu_k$ denote the distribution of the $k$-th iterate of SGLD-CV2 and let $\pi \propto e^{-U(\theta)}$ be the target posterior. let $\bar\theta$ be the mean of $\pi$. Then
\begin{equation}
    W_2^2(\nu_{k+1}, \pi) \leq \Big[(1-mh)\,W_2(\nu_k,\pi) + \alpha M(h^3 d)^{1/2}\Big]^2 + h^2\sigma^2_{\emph{CV2}},
\end{equation}
where $\alpha = 7\sqrt{2}/6$ and
\begin{equation}
    \sigma^2_{\emph{CV2}} = \frac{2N\sum_{i=1}^N M_i^2}{n}\left(8\sup_{k\geq 0}\,\mathbb{E}\big[\|\theta_k - \theta^*\|^4\big] + 8\|\theta^* - \bar\theta\|^4 + \|\hat\theta - \bar\theta\|^4\right),
\end{equation}
where $\theta^*$ is the global minimum of $U$. In particular, $W_2(\nu_k, \pi)$ contracts geometrically at rate $(1-mh)$, with bias term scaling asymptotically as $O(h\sigma_{\emph{CV2}}) = O(h\sqrt{N/n})$.}

\begin{proof}
Let $\pi$ be the invariant distribution of the underlying dynamics, so that it has density $e^{-U(\theta)} \propto p(\theta|\mathbf{x})$, and define $W_2(\nu_k, \pi)$ to be the Wasserstein distance between $\nu_k$ and $\pi$. Recall $\xi_k$ is the SGLD-CV2 gradient noise term, i.e. $\xi_k := \nabla \hat{U}(\theta_k) - \nabla U(\theta_k)$. Then we can write a single step of SGLD-CV2 as
\begin{equation}\label{eq:update-rule}
    \theta_{k+1} = \theta_k - h\nabla U(\theta_k) - h\xi_k + \sqrt{2h}\zeta_k,
\end{equation}

We have that $\theta_k \sim \nu_k$, and follow similarly to the proof of \citet{dalalyan2017}[Proposition 2]. First define $Y_0$ to be a draw from the invariant distribution $\pi$, such that the joint distribution of $Y_0$ and $\theta_k$ minimizes $\mathbb{E}\|Y_0 - \theta_k\|^2$. Here $\|\cdot\|$ denotes the Euclidean distance for $\mathbb{R}^d$. It follows that $\mathbb{E}\|Y_0 - \theta_k\|^2 = W_2^2(\nu_k, \pi)$.

Let $B_t$ be a $d$-dimensional Wiener process, independent of $\theta_k, Y_0$ and $\xi_k$ but which we couple to the injected noise $\zeta_k$ so that $B_h = \sqrt{h}\zeta_k$. Now let $Y_t, t > 0$, follow the diffusion
\begin{equation}
    Y_t = Y_0 - \int_0^t \nabla U(Y_s)ds + \sqrt{2}B_t.
\end{equation}

Let $\Delta_k = Y_0 - \theta_k$ and $\Delta_{k+1} = Y_h - \theta_{k+1}$. Since we started the process $Y_t$ from $Y_0 \sim \pi$, then it follows that $Y_t \sim \pi$ for all $t > 0$. Also since $W_2^2(\nu_{k+1}, \pi)$ minimizes the expected squared distance between two random variables with marginals $\nu_{k+1}$ and $\pi$ then it follows that
\begin{equation}
    W_2^2(\nu_{k+1}, \pi) \leq \mathbb{E}\|\Delta_{k+1}\|^2 \tag{$\star$}\label{eq:wass-id}
\end{equation}

Let us define
\begin{align}
    A &= \nabla U(\theta_k + \Delta_k) - \nabla U(\theta_k), \\
    V &= \int_0^h [\nabla U(Y_t) - \nabla U(Y_0)]dt.
\end{align}

Then by the unbiasedness of the gradient estimation, $\xi_k$ has mean 0 regardless of the value of $\theta_k$. Thus
\begin{align}
    \mathbb{E}\|\Delta_{k+1}\|^2 &= \mathbb{E}\|\Delta_k - hA - V\|^2 + h^2\mathbb{E}\|\xi_k\|^2 \notag \\
    &\leq \mathbb{E}\big[\|\Delta_k - hA\| + \|V\|\big]^2 + h^2\mathbb{E}\|\xi_k\|^2.
\end{align}

We can then apply Lemmas 2 and 4 in \citet{dalalyan2017}, stated below in Lemmas 4 and 5, as well as applying the gradient noise bound in Proposition 1, to obtain a bound on $W_2^2(\nu_{k+1}, \pi)$ given $W_2^2(\nu_k, \pi)$.   

\textbf{Lemma 4.} \textit{With $A$ as defined in (27), if $h < 2m/(2M^2 + m^2)$, then $\|\Delta_k - hA\| \leq (1 - mh)\|\Delta_k\|.$}

The original lemma by Dalalyan and Karagulyan (2017) assumed $h < 2/(m + M)$, but this holds when $h < 2m/(2M^2 + m^2)$ as $m \leq M$.

\textbf{Lemma 5.} \textit{Under Assumption 1. Let $V$ be as defined in (28), then}
\begin{equation}
    \mathbb{E}\|V\| \leq \frac{1}{2}(h^4 M^3 d)^{\frac{1}{2}} + \frac{2}{3}(2h^3 d)^{\frac{1}{2}}M.
\end{equation}

 Recall Proposition $\labelcref{prop:quartic}$:
\begin{align}
    \mathbb{E}\|\xi_k\|^2 &\leq \frac{N \sum_{i=1}^N M_i^2}{4n} \mathbb{E}\|\theta_k - \hat\theta\|^4 \notag
\end{align}
% Now consider the second expectation term on the RHS. Recall that $\bar\theta$ be the mean of the target posterior $\pi$:
% \begin{equation}\label{eq:sgldcv2var}
%     \mathbb{E}\|Y_0 - \hat\theta\|^4 \leq 8\mathbb{E}\|\hat\theta - \bar\theta\|^4 + 8\mathbb{E}\|Y_0 - \bar\theta\|^4.
% \end{equation}
% % Let $\theta^*$ denote the global minimum of the potential (the MAP estimate), such that $\nabla U(\bar\theta) = 0$.
% % TODO: check necessity
% It follows that
% \begin{equation}\label{eq:bound_4th}
%     \mathbb{E}\|\xi_k\|^2 \leq \frac{2N \sum M_i^2}{n} \mathbb{E}\|\theta_k - Y_0\|^4 + \frac{16N \sum M_i^2}{n} \left[ \mathbb{E}\|\hat\theta - \bar\theta\|^4 + \mathbb{E}\|Y_0 - \bar\theta\|^4 \right].
% \end{equation}

% Recall that $Y_0$ is drawn from the stationary distribution $\pi \propto \exp(-U(\theta))$. Under Assumption 1, the potential $U$ is $m$-strongly convex and the measure is $m$-strongly log-concave. By the classical Bakry-\'Emery criterion, strongly log-concave measures satisfy a Logarithmic Sobolev Inequality with a constant bounded by $1/m$ \citet{bakry_emery_2006}. When a probability measure satisfies a log-Sobolev inequality, all of its marginals exhibit sub-Gaussian tails. So $\pi$ is sub-Gaussian with a variance proxy of $1/m$, therefore its fourth moment scales as:
% \begin{equation}\label{eq:target_4th}
%     \mathbb{E}_{Y_0 \sim \pi}\|Y_0 - \bar\theta\|^4 \leq C_\pi \frac{d^2}{m^2},
% \end{equation}

% for some constant $C_\pi$ dependent on $\pi$. Now note that the second expectation term in $\labelcref{eq:bound_4th}$ is bounded by how good our found $\hat\theta$ in the Pre-processing step of the SGLD-CV2 algorithm is, i.e. how close it is to the posterior mean, so there is no further work to be done here to bound it.

We show that the fourth moment of the iterates, $\mathbb{E}\|\theta_k - \hat\theta\|^4$ is uniformly bounded for all $k \geq 0$ using a Foster-Lyapunov drift condition. To establish the restorative drift, we must center our analysis on the MAP estimate $\theta^*$ (where the true gradient $\nabla U(\theta^*) = 0$). Let $e_k = \theta_k - \theta^*$ and define the Lyapunov function:
$$V(\theta_k) = \|e_k\|^4$$

Write our SGLD update rule as:
$$e_{k+1} = e_k - hg_k + \sqrt{2h}\zeta_k \quad \quad \quad \quad g_k \coloneqq \nabla\hat{U}_{\text{CV2}}(\theta_k)$$

We expand:
$$\|e_{k+1}\|^4 = (\|e_k - hg_k\|^2 + 2\sqrt{2h}\langle e_k - hg_k, \zeta_k \rangle + 2h\|\zeta_k\|^2)^2$$

First we fix $g_k$ and take the expectation with respect to $\zeta$. Every term containing an odd power of $\zeta_k$ vanishes because the odd moments of a Gaussian are zero. Using the standard isotropic Gaussian properties of $\zeta_k \sim \mathcal{N}(0, I_d)$ we have:

\[
  \mathbb{E}_{\zeta}\bigl[\langle e_k - hg_k, \zeta_k \rangle^2\bigr] = \|e_k - hg_k\|^2,
  \qquad
  \mathbb{E}_{\zeta}\bigl[\|\zeta_k\|^2\bigr] = d,
  \qquad
  \mathbb{E}_{\zeta}\bigl[\|\zeta_k\|^4\bigr] = d(d+2).
\]

Then:
\begin{align}
\mathbb{E}_{\zeta}\bigl[\|e_{k+1}\|^4\bigr]
  &= \|e_k - hg_k\|^4
   + 4h(d+2)\|e_k - hg_k\|^2
   + 4h^2 d(d+2).
  \label{eq:after-zeta}
\end{align}

Let $\mathcal{F}_k$ be the sigma-field generated by the past up to time $k$.
We now take the final desired conditional expectation $\mathbb{E}_k[\,\cdot\,] \coloneqq
\mathbb{E}[\,\cdot \mid \mathcal{F}_k]$ of~\eqref{eq:after-zeta}.
Since $\zeta_k \sim \mathcal{N}(0, I_d)$ is independent of $\mathcal{F}_k$ and
the gradient estimator is unbiased, i.e.\
$\mathbb{E}[g_k \mid \mathcal{F}_k] = \nabla U(\theta_k)$:
\begin{equation}
  \|e_k - hg_k\|^2
  = \|e_k\|^2 - 2h\langle e_k, g_k \rangle + h^2\|g_k\|^2,
\end{equation}
so
\[
  \mathbb{E}_k\bigl[\|e_k - hg_k\|^2\bigr]
  = \|e_k\|^2
    - 2h\langle e_k, \nabla U(\theta_k) \rangle
    + h^2\,\mathbb{E}_k\bigl[\|g_k\|^2\bigr].
\]
\begin{align}
    \|e_k - hg_k\|^4 &= \|e_k\|^4 - 4h\|e_k\|^2\langle e_k, g_k \rangle + R_1, \\
    4h(d+2)\|e_k - hg_k\|^2 &= 4h(d+2)\|e_k\|^2 + R_2,
\end{align}
where the remainders $R_1$ and $R_2$ collect all higher-order terms in $h$:
\begin{align*}
    R_1 + R_2 &= 2h^2\|e_k\|^2\|g_k\|^2 + 4h^2\langle e_k, g_k \rangle^2 - 4h^3\langle e_k, g_k \rangle\|g_k\|^2 \\
    &\quad + h^4\|g_k\|^4 - 8h^2(d+2)\langle e_k, g_k \rangle + 4h^3(d+2)\|g_k\|^2.
\end{align*}

Under Assumption 1, the potential $U$ has a Lipschitz continuous gradient and Hessian. Consequently, the stochastic gradient $g_k$ has conditionally bounded moments such that $\mathbb{E}_k[\|g_k\|^2] \leq C_2(\|e_k\|^2 + 1)$ and $\mathbb{E}_k[\|g_k\|^4] \leq C_4(\|e_k\|^4 + 1)$ for some finite constants $C_2, C_4 > 0$. By applying the Cauchy-Schwarz inequality and Young's inequality ($\|e_k\|^2 \leq \|e_k\|^4 + 1$), the conditional expectation of these remaining terms can be upper bounded for all $h \leq 1$:
\begin{equation}
    \mathbb{E}_k[R_1 + R_2] \leq A h^2 \|e_k\|^4 + B h^2,
\end{equation}
for some positive constants $A, B > 0$ independent of $h$ and $e_k$. Taking $\mathbb{E}_k$ of \eqref{eq:after-zeta} and substituting our bounded remainder:
\begin{equation}\label{eq:drift_leading}
\begin{aligned}
    \mathbb{E}_k\big[\|e_{k+1}\|^4\big] &\leq \|e_k\|^4 - 4h\|e_k\|^2\langle e_k, \nabla U(\theta_k) \rangle \\
&\quad + 4h(d+2)\|e_k\|^2 + Ah^2\|e_k\|^4 + Bh^2 + 4h^2d(d+2).
\end{aligned}
\end{equation}

By Assumption 1, the potential $U$ is $m$-strongly convex. Since $\nabla U(\theta^*) = 0$, this guarantees the restorative drift:
\begin{equation}
    \langle \nabla U(\theta_k), e_k \rangle = \langle \nabla U(\theta_k) - \nabla U(\theta^*), \theta_k - \theta^* \rangle \geq m\|e_k\|^2.
\end{equation}
Substituting this into \eqref{eq:drift_leading} provides a negative leading drift of $-4mh\|e_k\|^4$. To handle the noise term $4h(d+2)\|e_k\|^2$ injected by the Brownian motion, we apply Young's inequality ($ab \leq \frac{\epsilon}{2}a^2 + \frac{1}{2\epsilon}b^2$). For any chosen $\epsilon > 0$, we have:
\begin{equation}
    4h(d+2)\|e_k\|^2 \leq h \left( \epsilon\|e_k\|^4 + \frac{4(d+2)^2}{\epsilon} \right).
\end{equation}

And so consolidating the constant terms $B' \coloneqq B + 4d(d+2)$ we have:
\begin{align}
    \mathbb{E}_k\big[\|e_{k+1}\|^4\big] &\leq \|e_k\|^4 - 4mh\|e_k\|^4 + \epsilon h\|e_k\|^4 + \frac{4h(d+2)^2}{\epsilon} + Ah^2\|e_k\|^4 + B'h^2 \notag\\
    &= \Big(1 - (4m - \epsilon)h + Ah^2\Big)\|e_k\|^4 + h\left( \frac{4(d+2)^2}{\epsilon} + B'h \right).
\end{align}

To show the Foster-Lyapunov drift condition is satisfied we wish to establish a geometric contraction by a factor of $(1 - \lambda h)$. Let $\lambda$ be any arbitrary constant in $(0, 4m)$. We can choose $\epsilon > 0$ sufficiently small such that $4m - \epsilon > \lambda$, which gives a positive gap $\delta \coloneqq 4m - \epsilon - \lambda > 0$. 

By the Archimedean principle, there exists a positive stepsize threshold $h_{\max} \coloneqq \min(1, \delta / A)$ such that for all $h \in (0, h_{\max}]$, we have $Ah \leq \delta$. Multiplying by $h$:
\begin{equation}
    Ah^2 \leq \delta h = (4m - \epsilon - \lambda)h.
\end{equation}
Substituting this upper bound into our coefficient for the fourth moment:
\begin{equation}
    1 - (4m - \epsilon)h + Ah^2 \leq 1 - (4m - \epsilon)h + (4m - \epsilon - \lambda)h = 1 - \lambda h.
\end{equation}

Let $$C_h \coloneqq \frac{4(d+2)^2}{\epsilon} + B'h_{\max}$$
we have a valid Foster-Lyapunov drift condition for any contraction rate $\lambda \in (0, 4m)$ and stepsizes $h \leq h_{\max}$:
\begin{equation}
    \mathbb{E}_k\big[\|\theta_{k+1} - \theta^*\|^4\big] \leq (1 - \lambda h)\|\theta_k - \theta^*\|^4 + C_h h.
\end{equation}

Taking the expectation and by Tower property:
\begin{equation}
     \mathbb{E}\big[\mathbb{E}_k\big[\|\theta_{k+1} - \theta^*\|^4\big]\big] \leq (1 - \lambda h)\mathbb{E}\|\theta_k - \theta^*\|^4 + \mathbb{E}[C_h h].
\end{equation}
Note that LHS$=\mathbb{E}\big[\mathbb{E}\big[\|\theta_{k+1} - \theta^*\|^4 \mid \mathcal{F}_k\big]\big]$ by definition of $\mathbb{E}_k$, and so by Tower property $\text{LHS}=\mathbb{E}\big[\|\theta_{k+1} - \theta^*\|^4 \big]$. So:
\begin{equation}
     \mathbb{E}\big[\|\theta_{k+1} - \theta^*\|^4 \big] \leq (1 - \lambda h)\mathbb{E}\|\theta_k - \theta^*\|^4 + \mathbb{E}[C_h h].
\end{equation}
Since $\lambda > 0$ and $C_h > 0$, by a discrete Gr\"onwall argument it follows that the sequence of fourth moments is uniformly bounded for all $k \geq 0$:
\begin{equation}
    B_4^* \coloneqq \sup_{k \geq 0} \mathbb{E}\big[\|\theta_k - \theta^*\|^4\big] < \infty.
\end{equation}

return to the variance bound of the SGLD-CV2 estimator from Lemma 1. We apply the triangle inequality once more to separate the iterates from the control variate centering value $\hat\theta$:
\begin{align}\label{eq:target-4th-bound}
    \mathbb{E}\|\xi_k\|^2 &\leq \frac{N \sum_{i=1}^N M_i^2}{4n} \mathbb{E}\|\theta_k - \hat\theta\|^4 \notag\\
    &\leq \frac{2N \sum_{i=1}^N M_i^2}{n} \left( \mathbb{E}\|\theta_k - \bar\theta\|^4 + \|\hat\theta - \bar\theta\|^4 \right)
    % &\leq \frac{2N \sum_{i=1}^N M_i^2}{n} \left( B_4 + \|\hat\theta - \bar\theta\|^4 \right) \equiv \sigma^2_{\text{CV2}}.
\end{align}

While the Foster-Lyapunov drift we've just established relies on centering at the MAP estimate $\theta^*$ (where $\nabla U(\theta^*) = 0$), in $\labelcref{eq:target-4th-bound}$, we wish to bound $\mathbb{E}\|\theta_k - \bar\theta\|^4$ which is evaluated relative to the posterior mean $\bar\theta$. For an $m$-strongly log-concave target posterior distribution, the distance between the MAP estimate and the true posterior mean $\|\theta^* - \bar\theta\|$ is deterministically bounded by a constant depending on the geometry of the posterior, thus there is no more work for us to do to bound this term. Applying inequality $\|a - b\|^4 \leq 8\|a - c\|^4 + 8\|b - c\|^4$ to $\mathbb{E}\|\theta_k - \bar\theta\|^4$, we get:
\begin{equation}
    \mathbb{E}\big[\|\theta_k - \bar\theta\|^4\big] \leq 8\mathbb{E}\big[\|\theta_k - \theta^*\|^4\big] + 8\|\theta^* - \bar\theta\|^4 \leq 8 B_4^* + 8\|\theta^* - \bar\theta\|^4 \equiv B_4.
\end{equation}

With the fourth moment of the iterates around the mean uniformly bounded by $B_4$, we return to the variance bound of the SGLD-CV2 estimator from $\labelcref{eq:target-4th-bound}$:
\begin{align}
    \mathbb{E}\|\xi_k\|^2 &\leq \frac{2N \sum_{i=1}^N M_i^2}{n} \left( \mathbb{E}\|\theta_k - \bar\theta\|^4 + \|\hat\theta - \bar\theta\|^4 \right) \notag\\
    &\leq \frac{2N \sum_{i=1}^N M_i^2}{n} \left( B_4 + \|\hat\theta - \bar\theta\|^4 \right) \equiv \sigma^2_{\text{CV2}}.
\end{align}

Note that $\sigma^2_{\text{CV2}}$ is a uniform global constant. Using $W_2^2(\nu_{k+1}, \pi) \leq \mathbb{E}\|\Delta_{k+1}\|^2$:
\begin{equation}
    W_2^2(\nu_{k+1}, \pi) \leq \mathbb{E}\big[\|\Delta_k - hA\| + \|V\|\big]^2 + h^2\mathbb{E}\|\xi_k\|^2.
\end{equation}
Applying the deterministic bounds from Lemmas 4 and 5, and substituting our uniform noise bound $\sigma^2_{\text{CV2}}$, we get the final recursive inequality:
\begin{equation}\label{eq:wass-contraction}
    W_2^2(\nu_{k+1}, \pi) \leq \Big[ (1 - mh)W_2(\nu_k, \pi) + \alpha M(h^3 d)^{\frac{1}{2}} \Big]^2 + h^2 \sigma^2_{\text{CV2}},
\end{equation}
where $\alpha = 7\sqrt{2}/6$. 

Note that noise variance $h^2 \sigma^2_{\text{CV2}}$ is an additive constant independent of $W_2(\nu_k, \pi)$, so our recursion will give a standard contractive sequence.
% Now using that $W_2^2(\nu_{k+1}, \pi) \leq \mathbb{E}\|\Delta_{k+1}\|^2$ we get the following
% \begin{align}
%     W_2^2(\nu_{k+1}, \pi) &\leq \left[ (1 - mh)W_2(\nu_k, \pi) + \alpha M(h^3 d)^{\frac{1}{2}} \right]^2 \notag\\
%     &\quad + \frac{2h^2 N \sum M_i^2}{n} \mathbb{E}\|\theta_k - Y_0\|^4 \notag\\
%     &\quad + \frac{16h^2 N \sum M_i^2}{n} \left[ \|\hat\theta - \bar\theta\|^4 + \mathbb{E}\|Y_0 - \bar\theta\|^4 \right],
% \end{align}
% where $\alpha = 7\sqrt{2}/6$. Gathering like terms we can further bound $W_2^2(\nu_{k+1}, \pi)$ to get a recursive 
\end{proof}
% \jay{we might want to use \texttt{\textbackslash Cref\{lemma:somelabel\}}}

\textbf{Theorem 2.} \legacytheoremlabel{2}{thm:sgld_cv2_convergence} \textit{Under Assumptions~\ref{ass:convex} and ~\ref{ass:grad}, let $\theta_K$ be the state of SGLD-CV2 at the $K^{th}$ iteration of the algorithm with stepsize $h$, initial value $\theta_0$, and control-variate centering value $\hat\theta$. Let the distribution of $\theta_K$ be $\nu_K$. Define $B_4 \coloneqq \sup_{k \geq 0} \mathbb{E}\big[\|\theta_k - \theta^*\|^4\big] + 8\|\theta^* - \bar\theta\|^4< \infty$. Define $R \coloneqq M/m$ as the condition number of the potential $U$. Then for any tolerance $\epsilon_0 > 0$ and any choice of error decomposition constants $c_1, c_2, c_3 > 0$ such that $c_1 + c_2 + c_3 = 1$, if the following conditions hold:}
\begin{align}
    h &\leq \frac{c_2^2 \epsilon_0^2}{R^2 \alpha^2 d m}, \label{eq:cond_h} \\
    K h &\geq \frac{1}{2m} \log\left[ \frac{m}{c_1^2 \epsilon_0^2} \left( \mathbb{E}\|\theta_0 - \bar\theta\|^2 + \frac{d}{m} \right) \right], \label{eq:cond_K} \\
    n &\geq \frac{2 N \sum_{i=1}^N M_i^2}{c_3^2 \epsilon_0^2 m} \left[ \|\hat\theta - \bar\theta\|^4 + B_4 \right], \label{eq:cond_n}
\end{align}
\textit{where $\alpha = 7\sqrt{2}/6$, and $d$ is the dimension of $\theta_k$, then the Wasserstein-2 distance to the target distribution satisfies $W_2(\nu_K, \pi) \leq \epsilon_0 / \sqrt{m}$.}

\begin{proof}
    Taking the square root of both sides of the Wasserstein contraction bound established for a single step of SGLD-CV2 from $\labelcref{eq:wass-contraction}$, and using that $\sqrt{x^2 + y^2} \leq x + y$ for $x,y$ positive (subadditivity of the square root function for positive numbers), we simplify the bound to:
\begin{equation}
    W_2(\nu_{k+1}, \pi) \leq (1-mh)W_2(\nu_k,\pi) + \alpha M h^{3/2} d^{1/2} + h\sigma_{\text{CV2}}.
\end{equation}

This recursive inequality over $K$ iterations gives:
\begin{align}
    W_2(\nu_K, \pi) &\leq (1-mh)^K W_2(\nu_0, \pi) + \sum_{j=0}^{K-1} (1-mh)^j \Big( \alpha M h^{3/2} d^{1/2} + h\sigma_{\text{CV2}} \Big) \notag \\
    &\leq (1-mh)^K W_2(\nu_0, \pi) + \frac{\alpha M h^{3/2} d^{1/2} + h\sigma_{\text{CV2}}}{1 - (1-mh)} \notag \\
    &= (1-mh)^K W_2(\nu_0, \pi) + \alpha \frac{M}{m} \sqrt{hd} + \frac{\sigma_{\text{CV2}}}{m}. \label{eq:unrolled_w2}
\end{align}

Suppose we stop the algorithm at iteration $K$ and require the total distance to be bounded by $W_2(\nu_K, \pi) \leq \epsilon_0 / \sqrt{m}$. Using \eqref{eq:unrolled_w2}, we can decompose this total error tolerance among the three terms using coefficients $c_1, c_2, c_3 > 0$ where $c_1 + c_2 + c_3 = 1$. A sufficient condition is to bound each error term respectively:
\begin{align}
    (1-mh)^K W_2(\nu_0, \pi) &\leq c_1 \frac{\epsilon_0}{\sqrt{m}}, \label{eq:bound_init} \\
    \alpha \frac{M}{m} \sqrt{hd} &\leq c_2 \frac{\epsilon_0}{\sqrt{m}}, \label{eq:bound_disc} \\
    \frac{\sigma_{\text{CV2}}}{m} &\leq c_3 \frac{\epsilon_0}{\sqrt{m}}. \label{eq:bound_noise}
\end{align}

First, we solve for the stepsize $h$ using the continuous-time discretization error bound \eqref{eq:bound_disc}. Recall that $R = M/m$:
\begin{equation}
    \alpha R \sqrt{hd} \leq c_2 \frac{\epsilon_0}{\sqrt{m}} \implies \alpha^2 R^2 hd \leq \frac{c_2^2 \epsilon_0^2}{m} \implies h \leq \frac{c_2^2 \epsilon_0^2}{R^2 \alpha^2 d m}.
\end{equation}
This is thus condition \eqref{eq:cond_h}.

Next, we solve for the required iterations $K$ using the initialization error bound \eqref{eq:bound_init}. Using the inequality $(1-mh)^K \leq \exp(-Kmh)$, we require:
\begin{equation}
    \exp(-Kmh) W_2(\nu_0, \pi) \leq c_1 \frac{\epsilon_0}{\sqrt{m}} \implies Kmh \geq \log\left( \frac{\sqrt{m} W_2(\nu_0, \pi)}{c_1 \epsilon_0} \right).
\end{equation}
Multiplying the argument of the logarithm by a factor of 2 (which squares the inside), we get:
\begin{equation}
    Kh \geq \frac{1}{2m} \log\left( \frac{m W_2^2(\nu_0, \pi)}{c_1^2 \epsilon_0^2} \right).
\end{equation}
Since the starting point $\theta_0$ is deterministic, we invoke Theorem 1 \cite{durmus2016highdimensional} which states $W_2^2(\nu_0, \pi) \leq \mathbb{E}\|\theta_0 - \bar\theta\|^2 + d/m$. Substituting this gives us condition \eqref{eq:cond_K}:
\begin{equation}
    Kh \geq \frac{1}{2m} \log\left[ \frac{m}{c_1^2 \epsilon_0^2} \left( \mathbb{E}\|\theta_0 - \bar\theta\|^2 + \frac{d}{m} \right) \right].
\end{equation}

Finally, we solve for the required batch size $n$ using \eqref{eq:bound_noise}:
\begin{equation}
    \frac{\sigma_{\text{CV2}}^2}{m^2} \leq \frac{c_3^2 \epsilon_0^2}{m} \implies \sigma_{\text{CV2}}^2 \leq c_3^2 m \epsilon_0^2.
\end{equation}
Substitute the definition of $\sigma^2_{\text{CV2}} = \frac{2N\sum M_i^2}{n}\big(\|\hat\theta - \bar\theta\|^4 + B_4\big)$:
\begin{equation}
    \frac{2N\sum_{i=1}^N M_i^2}{n}\Big(\|\hat\theta - \bar\theta\|^4 + B_4\Big) \leq c_3^2 m \epsilon_0^2.
\end{equation}
Rearranging for $n$ gives us condition \eqref{eq:cond_n}:
\begin{equation}
    n \geq \frac{2 N \sum_{i=1}^N M_i^2}{c_3^2 \epsilon_0^2 m} \Big( \|\hat\theta - \bar\theta\|^4 + B_4 \Big).
\end{equation}
Since all three sufficient conditions are met, $W_2(\nu_K, \pi) \leq \epsilon_0 / \sqrt{m}$ holds.
\end{proof}

\textbf{Remark.} If we allocate the error tolerance equally across all three sources of error symmetrically $c_1 = c_2 = c_3 = 1/3$, the multipliers evaluate to $c_2^2 = 1/9$, $1/c_1^2 = 9$, and $2/c_3^2 = 18$. Substituting we get the bounds:
\begin{align}
    h &\leq \frac{\epsilon_0^2}{9 R^2 \alpha^2 d m}, \\
    Kh &\geq \frac{1}{2m} \log\left[ \frac{9m}{\epsilon_0^2} \left( \mathbb{E}\|\theta_0 - \bar\theta\|^2 + \frac{d}{m} \right) \right], \\
    n &\geq \frac{18 N \sum_{i=1}^N M_i^2}{\epsilon_0^2 m} \left[ \|\hat\theta - \bar\theta\|^4 + B_4 \right].
\end{align}

\textbf{Corollary 1.} \textit{Assume that Assumptions~\ref{ass:convex} and ~\ref{ass:grad} and the conditions of Theorem 2 hold. Fix $\epsilon_0 > 0$ and let constants $c_1 = c_2 = c_3 = 1/3$. Define the cost constants:}
\begin{equation}
    C_1 = \frac{9 R^2 \alpha^2 d}{\epsilon_0^2}, \quad \text{and} \quad C_2 = \frac{18}{\epsilon_0^2}.
\end{equation}
\textit{We can implement an SGLD-CV2 algorithm achieving $W_2(\nu_K, \pi) \leq \epsilon_0/\sqrt{m}$ such that the computational cost $Kn$ is bounded by:}
\begin{equation}
    Kn \leq \left[ \frac{C_1}{2} \log\left( \frac{9m}{\epsilon_0^2}\left(\mathbb{E}\|\theta_0 - \bar\theta\|^2 + \frac{d}{m}\right) \right) \right] \left[ C_2 \frac{N \sum_{i=1}^N M_i^2}{m} \Big( \|\hat\theta - \bar\theta\|^4 + B_4 \Big) \right].
\end{equation}

\begin{proof}
To bound the total computational cost $Kn$, we assume the algorithm is run with the maximum allowed stepsize $h$ that satisfies condition \eqref{eq:cond_h} under symmetric allocations ($c_i=1/3$). By defining $C_1 = 9 R^2 \alpha^2 d / \epsilon_0^2$, we can write the maximum stepsize as:
\begin{equation}
    h_{\max} = \frac{1}{m C_1}.
\end{equation}
Substituting $h = h_{\max}$ and $c_1=1/3$ into the iteration required bound from condition \eqref{eq:cond_K}:
\begin{equation}
    K \geq \frac{1}{h_{\max}} \frac{1}{2m} \log\left[ \frac{9m}{\epsilon_0^2} \left( \mathbb{E}\|\theta_0 - \bar\theta\|^2 + \frac{d}{m} \right) \right] = \frac{C_1}{2} \log\left[ \frac{9m}{\epsilon_0^2} \left( \mathbb{E}\|\theta_0 - \bar\theta\|^2 + \frac{d}{m} \right) \right].
\end{equation}
Next, we rewrite the batch size bound from condition \eqref{eq:cond_n} using $c_3=1/3$ and the definition $C_2 = 18 / \epsilon_0^2$:
\begin{equation}
    n \geq C_2 \frac{N \sum_{i=1}^N M_i^2}{m} \Big( \|\hat\theta - \bar\theta\|^4 + B_4 \Big).
\end{equation}
The total computational cost to achieve the desired accuracy is bounded by multiplying the minimum required iterations $K$ by the minimum required batch size $n$. So:
\begin{equation}
    Kn \leq \left[ \frac{C_1}{2} \log\left( \frac{9m}{\epsilon_0^2}\left(\mathbb{E}\|\theta_0 - \bar\theta\|^2 + \frac{d}{m}\right) \right) \right] \left[ C_2 \frac{N \sum_{i=1}^N M_i^2}{m} \Big( \|\hat\theta - \bar\theta\|^4 + B_4 \Big) \right].
\end{equation}
\end{proof}

\textbf{Remark on Asymptotic Efficiency:} Corollary 1 implies a theoretical advantage of the second-order control variate in asymptotics over the first-order. The constants $C_1$ and $C_2$ depend on $\epsilon_0$ and the condition number $R = M/m$. As stated in Corollary 1 \cite{baker2017} $R$ is a constant as $N$ increases. Furthermore, under Assumption~\ref{ass:hess-unif} the sum of the squared Hessian bounds $\sum_{i=1}^N M_i^2$ increases linearly with $N$, making the factor $\frac{N \sum M_i^2}{m}$ scale as $\mathcal{O}(N)$. By the Bernstein-von Mises theorem (or Theorem 1 \cite{durmus2016highdimensional}), the variance of the posterior scales as $1/N$ as N increases, meaning the central fourth moment $B_4$ scales as $\mathcal{O}(N^{-2})$. This means that if in our preprocessing step we can find a mode $\hat\theta$ with SGD algorithm that is $\mathcal{O}(N^{-\frac{1}{2}})$ of the posterior mean $\bar\theta$, then $\left[ C_2 \frac{N \sum_{i=1}^N M_i^2}{m} \Big( \|\hat\theta - \bar\theta\|^4 + B_4 \Big) \right]$ will scales as $\mathcal{O}(N^{-1})$. From Corollary 1 \cite{baker2017}, the computational cost of SGLD-CV is the product of $2$ terms, in which the first term scales approximately the same as our first term, while the second term scales as $\mathcal{O}(m\|\hat\theta - \bar\theta\|^2)$ (where $m$ increases linearly with $N$ by Assumption~\ref{ass:convex}). Then the same mode we were discussing would give SGLD-CV a computational cost of $\mathcal{O}(1)$, while giving SLGD-CV2 a computational cost of $\mathcal{O}(1/N)$.

Because the algorithm requires $n \geq 1$, this indicates that as $N \to \infty$, the gradient noise of SGLD-CV2 decays so quickly that an $\mathcal{O}(1)$ mini-batch size is more than sufficient.

\subsection*{Setup Costs}
Like in Baker, we restate a result on the convergence of Stochastic Gradient Descent under the strongly convex conditions of Assumption~\ref{ass:convex} that would give us the cost of the preprocessing step to find a desired mode as discussed above. Adapted from Nemirovski et al. (2009):

\textbf{Lemma 3.} \textit{Under Assumption~\ref{ass:convex}, let $\hat\theta$ denote the final state of SGD with stepsizes $h_k = 1/(mk)$ after $K_{\text{setup}}$ iterations. Suppose $\mathbb{E}\|\nabla U(\theta)\|^2 \leq D^2$ and denote the true MAP estimate (the global minimum of $U$) by $\theta^*$. Then it holds that:}
\begin{equation}
    \mathbb{E}\|\hat\theta - \theta^*\|^2 \leq \frac{4D^2}{m^2 K_{\text{setup}}}.
\end{equation}

Again as stated in Lemma 3 \cite{baker2017}, if we assume that the norm of the individual gradients $\|\nabla U_i(\theta)\|$ is bounded by some constant $\sigma$ for all $i \in \{1, \ldots, N\}$, then $D^2$ will be $\mathcal{O}(N^2)$. Because $m$ scales as $\mathcal{O}(N)$, the ratio $\frac{D^2}{m^2}$ reduces to $\mathcal{O}(1)$. To ensure that our squared error $\mathbb{E}\|\hat\theta - \theta^*\|^2$ scales as $\mathcal{O}(N^{-1})$—which implies the necessary $\mathcal{O}(N^{-2})$ fourth-moment scaling for our algorithm—we need $K_{\text{setup}}$ to be $\mathcal{O}(N)$. This means that in our preprocessing step we need to process the full dataset for SGD to converge to the desired mode $\hat\theta$ as such.

It follows that, just as in the first-order case, there are two one-off $\mathcal{O}(N)$ setup costs: one to run SGD to find an acceptable mode $\hat\theta$ within $\mathcal{O}(N^{-1/2})$ of the true mode, and one to evaluate the full dataset once to compute the exact log-posterior gradient $\nabla U(\hat\theta)$ to initialize the control variate.
